\section{Synthetic spectrum calculation in \textsc{NebulaPy}}
\label{sec:nebulapy-spectrum}

\textsc{NebulaPy} constructs an intrinsic, optically thin spectrum directly from
the thermodynamic state and ion populations stored in a PION simulation.  The
calculation combines bound--bound line emission, free--free emission
(bremsstrahlung), free--bound emission (radiative recombination), and
two-photon emission.  The atomic structure and much of the atomic emissivity
physics are supplied by CHIANTI through \textsc{ChiantiPy}, whereas the mapping
from PION variables to CHIANTI ions, the treatment of arbitrary non-equilibrium
ion populations, the cell-by-cell normalisation, grid masking, volume
integration, multiprocessing, and final spectral assembly are implemented in
\textsc{NebulaPy}.  The result is a source luminosity spectrum per unit solid
angle,
\begin{equation}
  \frac{\mathrm{d}L_\lambda}{\mathrm{d}\Omega}
  \quad [\mathrm{erg\,s^{-1}\,sr^{-1}\,\mathring{A}^{-1}}] ,
\end{equation}
and is not a detector-folded count spectrum.  For isotropic emission, the
angle-integrated luminosity is obtained by multiplying by $4\pi$.

\subsection{Input plasma state and ion selection}

For each simulation snapshot, the driver script reads the temperature
$T_c$, mass density $\rho_c$, electron density $n_{\mathrm e,c}$, ion mass
fractions $X_{Z,q,c}$, cell volumes $V_c$, and native nested-grid mask.  Here
$c$ denotes a computational cell, $Z$ an element, and $q$ its ionic charge.
The ion number density used by the spectrum calculation is evaluated by
\textsc{NebulaPy} as
\begin{equation}
  n_{Z,q,c} = \frac{\rho_c X_{Z,q,c}}{m_Z},
  \label{eq:ion-number-density}
\end{equation}
where $m_Z$ is the atomic mass used by the code.  If the highest ionisation
stage is not stored as a separate tracer, its mass fraction is reconstructed
as the elemental mass fraction minus the sum of the explicitly followed lower
stages; small negative residuals are clipped to zero.  The electron density is
calculated consistently from the charge-weighted ion mixture, with a small
numerical floor applied where necessary.

PION ion labels such as \texttt{O6+} are translated by \textsc{NebulaPy} to
CHIANTI notation such as \texttt{o\_7}.  The \textsc{ChiantiPy}
\texttt{specTrails.ionGate} machinery is then used to discover which ions and
radiative processes are available in the CHIANTI database.  The user may
request a restricted ion list or all available ions belonging to the selected
elements.  In the latter case, the snapshot driver obtains all corresponding
ion-density fields from PION before spectral synthesis.

The emitting domain can be restricted without altering the underlying
simulation data.  In the example X-ray workflow, \textsc{NebulaPy} reloads the
native grid mask for every snapshot and sets it to zero outside
$10^6\leq T/\mathrm{K}\leq10^9$.  Defining the resulting multiplicative mask as
$M_c$, a value of zero excludes a cell, unity includes it completely, and a
fractional value acts as a filling factor.  Independently of this user filter,
the spectrum kernel accepts coefficients only for
$10^2\leq T/\mathrm{K}\leq10^9$.

\subsection{Wavelength grid}

The spectral interval may be given either as wavelength bounds or as photon
energy bounds.  Energy bounds are converted using
\begin{equation}
  \lambda\,[\mathring{\mathrm A}]
  = \frac{12.39841984}{E\,[\mathrm{keV}]},
\end{equation}
with the order of the endpoints reversed because energy and wavelength are
inversely related.  Two wavelength grids are supported.  A uniform user grid
contains a prescribed number of equally spaced points between the limits.  The
alternative grid is constructed by querying the CHIANTI transition wavelengths
of every selected line-emitting ion, retaining wavelengths inside the requested
interval, adding the exact endpoints, removing duplicates, and sorting the
result.  All ions and all emission mechanisms are evaluated on this common
grid.

\subsection{Local emissivity coefficients}

For compactness, let
$C^{(p)}_{Z,q}(\lambda;T_c,n_{\mathrm e,c})$ denote the coefficient returned
for process $p$.  \textsc{NebulaPy} deliberately prevents CHIANTI from applying
an elemental abundance or a collisional-ionisation-equilibrium (CIE) fraction;
these factors are instead represented by the physical ion density in
Equation~\eqref{eq:ion-number-density}.  The cell contribution for the
continuum and the density-normalised line coefficients can therefore be
written as
\begin{equation}
  \left.\frac{\mathrm dL_\lambda}{\mathrm d\Omega}\right|_c
  = M_c V_c n_{\mathrm e,c}
    \sum_{Z,q}\sum_p n_{Z,q,c}
    C^{(p)}_{Z,q}(\lambda;T_c,n_{\mathrm e,c}).
  \label{eq:cell-spectrum}
\end{equation}
The implementation forms the weight
$w_c=M_cV_cn_{\mathrm e,c}$, multiplies it by the relevant ion density, and
contracts the cell axis with the coefficient matrix using a vectorised tensor
operation.  This is mathematically equivalent to an emission-measure
integration but avoids constructing an intermediate, temperature-binned
differential emission measure.

\paragraph{Free--free emission.}
The continuum object and free--free Gaunt factors are obtained from
\textsc{ChiantiPy}.  Specifically, the Itoh Gaunt factor is used where defined
and the Sutherland value replaces any undefined Itoh entries.  The surrounding
coefficient calculation is implemented in \textsc{NebulaPy}.  For an ion of
charge $q$, its wavelength dependence has the form
\begin{equation}
  C^{(\mathrm{ff})}_{Z,q}(\lambda,T)
  = A_{\mathrm{ff}}\,
    \frac{q^2}{T^{1/2}\lambda^2}
    \exp\!\left(-\frac{hc}{\lambda k_{\mathrm B}T}\right)
    g_{\mathrm{ff}}(\lambda,T),
  \label{eq:free-free}
\end{equation}
where $A_{\mathrm{ff}}$ is assembled in \textsc{NebulaPy} from the physical
constants in the code.  This coefficient is multiplied by
$n_{\mathrm e}n_{Z,q}V$ in Equation~\eqref{eq:cell-spectrum}.

\paragraph{Free--bound emission.}
The free--bound coefficient is taken from
\textsc{ChiantiPy}'s \texttt{continuum.freeBound} calculation.  This uses the
CHIANTI free-bound level data, Verner photoionisation cross sections for the
ground-state contribution (when enabled), and Karzas--Latter Gaunt factors for
excited levels.  \textsc{NebulaPy} calls it with
\texttt{includeAbund=False}, \texttt{includeIoneq=False}, and no supplied
emission measure.  Thus neither a database abundance nor a CHIANTI CIE ion
fraction is folded into the returned spectrum.  In the current implementation,
the returned coefficient for a selected CHIANTI ion is weighted by the electron
density and the number density carrying the same ion label in the PION--CHIANTI
mapping.  This convention should be stated when comparing the result with codes
that label a recombination continuum by the recombining parent ion rather than
the target ion.

\paragraph{Line emission.}
Level populations and transition emissivities are evaluated by
\textsc{ChiantiPy}'s \texttt{ion.emiss} routine.  Its line power includes the
photon energy and $1/(4\pi)$ solid-angle factor and excludes elemental abundance
and ionisation fraction.  Because the CHIANTI line emissivity returned at this
stage contains the electron-density dependence used in solving the statistical
equilibrium populations, \textsc{NebulaPy} divides it by the supplied
$n_{\mathrm e}$ (subject to the numerical density floor) before applying the
common $n_{\mathrm e}n_{Z,q}V$ weight in
Equation~\eqref{eq:cell-spectrum}.  This prevents the electron density from
being counted twice.

Each discrete line of rest wavelength $\lambda_0$ and integrated power
$\epsilon_{ul}$ is distributed on the wavelength grid using a profile
$\phi(\lambda;\lambda_0)$ normalised so that
\begin{equation}
  \int\phi(\lambda;\lambda_0)\,\mathrm d\lambda=1,
  \qquad
  C^{(\mathrm{line})}(\lambda)
  =\sum_{u>l}\epsilon_{ul}\phi(\lambda;\lambda_{ul}).
\end{equation}
By default, the profile is \textsc{ChiantiPy}'s
\texttt{filters.gaussianR}; \textsc{NebulaPy} chooses a grid-aware resolving
factor, capped at $10^4$, so that the narrowest selected line is sampled by
approximately five grid intervals across its full width at half maximum.
Alternatively, the new \textsc{NebulaPy} global-profile implementation applies
a bulk Doppler shift and combines instrumental and velocity broadening.  For
the classical Doppler option,
\begin{align}
  \lambda_{\rm c} &= \lambda_0\left(1+\frac{v}{c}\right),\\
  \sigma_{\rm inst} &=
    \frac{\lambda_{\rm c}}{2\sqrt{2\ln2}\,R},\\
  \sigma_v &= \lambda_{\rm c}\frac{\sigma_{v,\rm los}}{c},\\
  \sigma_\lambda &=
    \left(\sigma_{\rm inst}^2+\sigma_v^2\right)^{1/2}.
\end{align}
The relativistic longitudinal Doppler factor may be selected instead.  The WR
140 example uses $v=-1365\,\mathrm{km\,s^{-1}}$,
$\sigma_{v,\rm los}=578\,\mathrm{km\,s^{-1}}$, and no instrumental broadening
($R=0$).

\paragraph{Two-photon emission.}
Two-photon emission is evaluated only for non-dielectronic H-like and He-like
ions for which the requested wavelength grid extends beyond the appropriate
transition threshold.  Atomic levels and the spectral emissivity are supplied
by \textsc{ChiantiPy}'s \texttt{ion.twoPhotonEmiss}.  Since that routine does
not use the same per-steradian normalisation as the other selected outputs,
\textsc{NebulaPy} divides its result by $4\pi$ before adding it to the spectrum.

\subsection{Non-equilibrium ionisation spectrum}

The distinction between an equilibrium and a non-equilibrium spectrum enters
through the ion populations, not through a different atomic emissivity solver.
For CIE synthesis one would write
\begin{equation}
  n^{\rm CIE}_{Z,q,c}
  = n_{Z,c}\,f^{\rm CIE}_{Z,q}(T_c),
\end{equation}
where the ion fraction $f^{\rm CIE}_{Z,q}$ is a single-valued function of the
local temperature (for a specified equilibrium model).  In the non-equilibrium
ionisation (NEI) workflow, \textsc{NebulaPy} instead reads the instantaneous,
advected ion mass fractions evolved by the PION/NEMO simulation and converts
them to $n^{\rm NEI}_{Z,q,c}$ using Equation~\eqref{eq:ion-number-density}.  It
then evaluates the same CHIANTI atomic coefficients at the local $T_c$ and
$n_{\mathrm e,c}$ but weights them by these simulation ion densities:
\begin{equation}
  \frac{\mathrm dL^{\rm NEI}_\lambda}{\mathrm d\Omega}
  = \sum_c M_cV_cn_{\mathrm e,c}
    \sum_{Z,q}n^{\rm NEI}_{Z,q,c}
    \sum_p C^{(p)}_{Z,q}(\lambda;T_c,n_{\mathrm e,c}).
  \label{eq:nei-spectrum}
\end{equation}
Consequently, two cells with the same present-day temperature and density may
emit different spectra if their thermal and ionisation histories have produced
different ion fractions.  This permits under-ionised gas behind a rapidly
heated shock and over-ionised gas in a rapidly cooling flow to be represented
without forcing the charge-state distribution back to CIE.  Importantly, the
\textsc{NebulaPy} spectrum module does \emph{not} itself integrate the
time-dependent ionisation and recombination rate equations.  That evolution is
performed by the simulation; spectral synthesis is a post-processing step that
preserves the resulting NEI populations.  Setting
\texttt{includeIoneq=False} in the CHIANTI continuum call and using
\texttt{ion.emiss}, which does not apply an ionisation fraction, are therefore
essential: applying the CHIANTI equilibrium fraction as well would erase or
double-weight the simulated non-equilibrium charge-state information.

\subsection{Parallel integration and output representation}

For a two-dimensional nested grid, one complete row at one refinement level is
submitted as a multiprocessing task.  Within a worker, CHIANTI evaluates all
cells of the row vectorially.  Each process coefficient is immediately
multiplied by $M_cV_cn_{\mathrm e,c}n_{Z,q,c}$ and summed over the row.  Only
one ion-by-wavelength vector per row is returned to the parent process, which
reduces inter-process communication and avoids retaining cell-resolved spectral
cubes.  The parent accumulates rows, then ions, to produce the one-dimensional
array in Equation~\eqref{eq:nei-spectrum}.

The native output is stored on the increasing wavelength grid as
$\mathrm dL_\lambda/\mathrm d\Omega$.  The example driver also supports energy
and photon-number representations.  With
$E=K/\lambda$, $K=12.39841984\,\mathrm{keV\,\mathring A}$, and
$1\,\mathrm{keV}=1.602176634\times10^{-9}\,\mathrm{erg}$,
\begin{align}
  \frac{\mathrm dL_E}{\mathrm d\Omega}
    &= \frac{\mathrm dL_\lambda}{\mathrm d\Omega}
       \left|\frac{\mathrm d\lambda}{\mathrm dE}\right|
     = \frac{\mathrm dL_\lambda}{\mathrm d\Omega}\frac{K}{E^2},\\
  \frac{\mathrm dN_E}{\mathrm d\Omega}
    &= \frac{1}{E\,(1.602176634\times10^{-9})}
       \frac{\mathrm dL_E}{\mathrm d\Omega}.
\end{align}
When plotted against energy the arrays are reordered so that the horizontal
axis increases, but this plotting transformation does not change the stored
wavelength-space spectrum.

\subsection{Division of responsibility between \textsc{ChiantiPy} and
\textsc{NebulaPy}}

In summary, the CHIANTI database and \textsc{ChiantiPy} supply the atomic level
and transition data, ion metadata and process availability, statistical-
equilibrium line populations and line powers, free-bound continuum physics,
free--free Gaunt factors, two-photon emissivities, and the default Gaussian
line-filter function.  \textsc{NebulaPy} supplies the PION-to-CHIANTI naming
interface; extraction and reconstruction of simulation ion densities; explicit
removal of CHIANTI abundance and CIE-fraction weighting; the free--free
coefficient assembly around the CHIANTI Gaunt factors; wavelength-grid
construction; line selection and summation; the optional bulk-shift and global
broadening model; cell masks, volumes, and temperature cuts; the NEI
normalisation by the simulation charge-state densities; row-wise parallel
reduction; summation over ions and processes; and conversion to the requested
plotting representation.
